%!%
% Copyright (c) 2026 mingcheng <mingcheng@outlook.com>
% 
% File: ch01-icebreaker.tex
% Author: mingcheng <mingcheng@outlook.com>
% File Created: 2026-03-27 14:57:13
% 
% Modified By: mingcheng <mingcheng@outlook.com>
% Last Modified: 2026-03-29 07:14:37
%%

% !TeX root = ../prompt-engineering-guide-for-beginners.tex

\chapter{破冰：把 AI 当成你的超级实习生}
\label{ch:icebreaker}

如果你觉得 AI 离自己很远，甚至有些高深莫测，那这一章就是专门为你准备的。
读完之后你会发现，AI 既不神秘，也不可怕。
它本质上就是一个特别能写字的助手，而关键在于你怎么跟它「说话」。

\section{什么是提示工程（Prompt Engineering）？}

简单来说，提示工程就是「面向大型概率模型进行自然语言编程」的学问。

你在 AI 对话框里输入的每一段文字，都叫做「提示词」（Prompt）。而提示工程研究的，就是如何写出更好的提示词，这实际上是在为模型设定一个精密的“初始上下文”，从而引导并在数学概率上极大限制它生成你想要的结果。你可以把它理解为新时代的「黑盒编程语言」：只不过你用的不再是需要编译和严谨语法的 Python 或 C++，而是我们日常交流、具有一定歧义的自然语言。

打个比方：去餐厅点菜时，如果只说「来个菜」，服务员完全不知道你想吃什么；但如果你说「来一份不辣的番茄炒蛋，少油少盐」，厨师就能精准地做出你想要的菜。跟 AI 打交道，道理完全一样。

提示工程不要求你会写代码，也不要求你懂机器学习算法。它是一项\textbf{人人都能学会的沟通技能}，只要你会打字、能表达需求，就可以上手。

\begin{tipbox}
很多新手在初次使用 AI 时，会担心自己的表达不够书面化或有语病，实际上大模型的语义理解能力非常强，对语言的容错率极高。就像与真人助理交代工作一样用大白话提问即可，完全不必过分纠结语法与辞藻。
\end{tipbox}

\section{大语言模型到底是个啥？}

你肯定听过 ChatGPT、Claude、Gemini 这些名字。它们背后的核心技术叫做「大语言模型」（Large Language Model，简称 LLM）。

不用被术语吓到。从技术底层来看，你可以把大语言模型想象成一个基于 Transformer 架构的「超级文字接龙机器」：它在训练阶段阅读过数十 TB 的乃至更庞大的互联网数据（包括文章、书籍和代码库），通过自注意力机制（Self-Attention）死记硬背下了人类语言中词与词之间的高维统计学规律。当你给它提供一句话时，它其实是在计算概率分布，预测接下来最可能出现的下一个词（Token），如此循环往复，一个词一个词地「接」出一段完整的回答。比如，当你输入「白日依山」，它就会自然而然地算出「尽」的出现概率最高（可能高达 99\%），从而接上这个字。

不仅如此，现代的对话大模型（如 ChatGPT）还经过了“指令微调（SFT）”和“人类反馈强化学习（RLHF）”。如果是原始的“基础模型”，它只会单纯接龙；但在经过指令微调后，它才真正变成了一个能听懂人类指令并“服从命令”的 AI 助手。

换句话说，AI 并不是真的像人类一样在利用物理大脑「理解」你说了什么，也不是在用独立意识「思考」问题的答案。它的本质是在做一个极其复杂的数学概率运算，生成一段在人类看来连贯且符合常理的响应。

这意味着两件事：

\begin{itemize}
  \item 它能写出非常流畅、专业的内容，因为它见过的好文章足够多；
  \item 它也可能一本正经地胡说八道，因为它只在意「说得通」，并不在意「是不是真的」。
\end{itemize}

理解这一点很重要，因为它直接决定了我们应该怎么使用 AI、怎么看待它给出的回答。

\section{为什么你的 AI 看起来很笨？}

很多人第一次用 AI 时，问了一个问题，得到一个不太满意的回答，就下结论说「AI 不行」。

但真相往往是：不是 AI 不行，而是我们提问的方式有问题。

来看一个对比：

\begin{itemize}
  \item \textbf{模糊的提问}：「帮我写个东西。」
  \item \textbf{清晰的提问}：「帮我写一封给客户的邮件，告诉对方项目进度延迟了两周，语气要礼貌但诚恳，控制在 200 字以内。」
\end{itemize}

第一种问法，AI 根本不知道你想要什么，只能瞎猜，给出的内容往往大而无当。第二种问法，AI 有了明确的上下文（特定的对象、背景、语气和字数限制），自然能给出靠谱得多的回答。

这就是提示工程中的一条黄金法则：\textbf{垃圾输入 = 垃圾输出}（Garbage In, Garbage Out）。你给 AI 的指令越模糊，回答就越空洞；指令越具体，回答就越精准。为了让 AI 表现得像个专业的得力助手，你首先得给它提供清晰、充分的背景信息。

再举一个日常例子。假设你想让 AI 帮你策划一顿朋友生日聚餐：

\begin{itemize}
  \item \textbf{模糊版}：「帮我推荐一家餐厅。」
  \item \textbf{清晰版}：「帮我推荐一家上海静安区适合 6 人聚餐的餐厅，预算人均 200 元左右，朋友中有一个人吃素，最好有独立包间，氛围适合过生日。」
\end{itemize}

前者，AI 可能给你推荐从街边麻辣烫到高端法餐的一串清单，什么都有但什么都没用。后者，AI 就能在地段、人数、预算、饮食限制等多项约束条件的交叉筛选下给出靠谱得多的答案。

好消息是，学会「好好提问」并不难，接下来的章节会一步步带你掌握。

\begin{tipbox}
如果 AI 给出的一次回答不尽如人意，不要立刻开启新对话或认为求助失败。你应该直接在对话框里追问和修正，例如补充一句「总结得太长了，帮我缩短到 100 字」或「语气再活泼一些」。这种不断微调的「多轮对话」，本身就是提示工程中非常高频且关键的操作。
\end{tipbox}

\section{市面上的模型怎么选？}

目前市面上有很多 AI 工具可供选择，各有特点。以下是一些主流选项的简要介绍：

\begin{itemize}
  \item \href{https://chatgpt.com}{\textbf{ChatGPT}}（OpenAI）：目前最知名的 AI 对话工具，功能全面，支持文本、图片、代码等多种任务，有免费版和付费版。
  \item \href{https://claude.ai}{\textbf{Claude}}（Anthropic）：擅长长文本处理和细致的分析任务，语气相对温和严谨，适合需要高质量文字输出的场景。
  \item \href{https://gemini.google.com}{\textbf{Gemini}}（Google）：谷歌推出的 AI 助手，和 Google 生态深度整合，适合习惯使用谷歌全家桶的用户。
  \item \href{https://qwen.ai}{\textbf{Qwen}}（阿里巴巴）: 阿里巴巴的 AI 模型，支持多语言输入，适合需要处理中文和其他语言混合内容的用户。
  \item \href{https://kimi.moonshot.cn}{\textbf{Kimi}}（月之暗面）：国产 AI 工具，支持超长文本输入，适合需要阅读和分析大量中文资料的用户。
  \item \href{https://www.deepseek.com}{\textbf{DeepSeek}}：国产开源模型的代表，性价比高，在推理和代码方面表现不俗。
  \item \href{https://www.zhipuai.cn}{\textbf{智谱 GLM}}：清华背景的国产大模型，提供丰富的 API 接口，适合有开发需求的用户。
\end{itemize}

如果你是纯新手，建议从 ChatGPT、Qwen、Kimi 或 DeepSeek 中任选一个开始。它们注册简单、界面友好，免费额度足够入门练习。不必纠结哪个「最好」，对于日常使用来说，\textbf{哪个顺手就用哪个}。随着你对 AI 了解的不断加深，自然会体会到不同模型在特定任务上的细微优势。

值得一提的是，目前很多平台都提供了「深度思考」或「推理模型」的选项（如 OpenAI 的 o1/o3、DeepSeek R1 等）。这类模型在回答前会先进行一段内部推理，适合处理数学、逻辑和复杂分析等任务。入门阶段用普通的对话模型就足够了，关于推理模型的更多细节将在第~\ref{sec:reasoning-models}~节介绍。

顺便推荐一下 \href{https://openrouter.ai}{OpenRouter}、\href{https://zenmux.ai}{ZenMux} 等聚合平台，它们提供统一的接口服务，让你可以在不同模型之间自由切换和对比试用，方便找到最适合自己的那个。
此外，\href{https://cherry-ai.com}{Cherry Studio}、\href{https://anythingllm.com}{AnythingLLM} 等第三方客户端工具也值得关注，它们不仅提供了多模型切换的便利，还内置了提示词管理功能，能帮助你更高效地组织与沉淀自己的提示词库，是进阶使用者的好帮手。

我个人使用 Claude 和 Gemini 比较多，但说实话 Claude 的费用偏高。如果你只是想练习提示词写作，完全可以先用 Qwen 或 Kimi 等有免费额度的模型，等熟练之后再考虑升级到更强大的付费模型。

\begin{tipbox}
无论选择哪个模型，在工作中使用 AI 时都请务必注意数据安全与隐私。请勿将公司的核心商业机密、内部代码库的密钥或客户的敏感个人信息毫无处理地直接输入给公共 AI 服务。
\end{tipbox}

\section{AI 的能力边界：它能做什么、不能做什么}

在深入学习之前，有必要先了解 AI 的能力边界，这样可以少走弯路。

\subsection*{AI 擅长的事情}

AI 最大的优势在于处理文字相关的任务：速度快、不怕重复、不会抱怨加班。具体来说，它擅长以下几类工作：

\begin{itemize}
  \item \textbf{文本生成}：写邮件、写文案、写报告、写故事。给一个方向，它就能快速产出初稿，省去「从零开始」的痛苦。
  \item \textbf{翻译}：中英互译乃至多语言翻译。日常沟通级别的翻译，AI 已经足够好用。
  \item \textbf{总结和提炼}：帮你把长文章浓缩成几句话。面对一份 50 页的报告，让 AI 先帮你抓重点，效率提升显著。
  \item \textbf{头脑风暴}：提供创意和灵感。当你卡在一个想法上不知如何展开时，AI 是个不错的「点子发生器」。
  \item \textbf{格式转换}：把文字整理成表格、列表、JSON 等格式。这类机械性的排版工作，交给 AI 再合适不过。
  \item \textbf{代码辅助}：写代码、改 Bug、解释代码逻辑。即使你不是程序员，AI 也能帮你写一些简单的自动化脚本。
\end{itemize}

\begin{tipbox}
遇到复杂任务时，不要指望 AI 能「一次写出完美成品」。最好的心态是把它当成一个不知疲倦的「草稿生成器」。由它快速完成从 0 到 80 分的草图与体力活，你再负责从 80 到 100 分的精修与灵魂注入，这才是目前最高效的人机协作模式。
\end{tipbox}

\subsection*{AI 不擅长的事情}

与此同时，AI 也有明显的短板。认清这些不足，能帮你避免在错误的场景中使用它：

\begin{itemize}
  \item \textbf{实时信息}：如果没有联网功能，AI 不知道今天的新闻或最新的股价。它的知识有一个「截止日期」，此后的事情一概不知。
  \item \textbf{精确计算}：复杂的数学运算它可能算错。AI 不是计算器，它是在「预测下一个最可能的词」，而不是在「算数」。虽然新版模型在这方面持续改进（例如能够自动调用代码来辅助计算），但涉及重要财务或业务数据时，仍建议自行验算。
  \item \textbf{事实准确性}：它可能会「编造」听起来很真实的假信息（这在一定程度上被称为「幻觉」，Hallucination）。比如你问它某本冷门书籍的作者，它可能自信满满地报出一个错误的名字。因此，在查阅专业文献或获取硬核知识时，务必进行交叉验证。
  \item \textbf{读心术}：你没说的东西，它猜不到。AI 只能根据你提供的信息来工作，无法感知你的真实意图、工作背景或情绪状态。
\end{itemize}

值得注意的是，虽然越来越多的模型已经支持联网搜索、调用外部工具等扩展功能（例如 MCP，即 Model Context Protocol，一种让 AI 连接外部数据源的标准化协议），但这并不意味着它就能完全克服上述短板。你仍然需要对 AI 的回答保持审慎态度，尤其是涉及重要决策或敏感信息时。

\subsection*{技术视角的常见误解}

最后，澄清几个对于 LLM 底层运作常见的「认知误区」：

\begin{itemize}
  \item \textbf{大模型不是搜索引擎}。搜索引擎是去互联网倒排索引中查找已有的映射网页，而 AI 则是根据学到的语言分布模式凭空“计算”并生成新的文字。它不是在特定的硬盘上「查」资料，而是在神经网络权重中「推演」回答。
  \item \textbf{大模型不是关系型数据库}。它没有内部的一个 SQL 乃至图数据库供你完美精确的提取信息。它的知识是模糊地压缩在千亿个神经元参数里的。因此绝对不能把它当成严格的词典或百科全书来信任。
  \item \textbf{无状态（Stateless）本质：AI 默认不具有“长久记忆”}。作为纯粹的推理引擎，每次你开启新对话，它的计算状态都会被清零。它就像一个「永远清醒但频繁失忆」的打字机，对你上个窗口说过的话一无所知（除非平台本身开发了额外的记忆检索外挂体系并自动组装到了你的系统提示词里）。
\end{itemize}

从合理的期待出发，逐步探索 AI 的能力边界，你会发现它是一个非常强大的生产力工具，能在很多场景下帮你节省大量时间和精力。
了解了这些能力边界，你就能把 AI 放在最擅长的位置上使用，既不会因期望过高而失望，也不会因误解而低估它的价值。

理论已经铺设完毕，接下来最好的学习方式就是亲自上手。下一章，我们会在五分钟内完成你的第一次 AI 对话。
